% Part: second-order-logic
% Chapter: metatheory
% Section: compactness

\documentclass[../../../include/open-logic-section]{subfiles}

\begin{document}

\olfileid{sol}{met}{lst}

\olsection{Löwenheim–Skolem定理对二阶逻辑不成立}

\begin{explain}
（向下）勒文海姆–斯科伦定理断言，每个有无限模型的语句集都有一个可数模型。
它也是完备性定理的一个推论：完备性证明为任何一致的语句集构造出一个模型，而那个模型是可数的。
还有一个向上勒文海姆–斯科伦定理，它保证如果一个语句集有一个可数无穷模型，那么它也有一个不可数模型。
这两个定理在二阶逻辑中都不成立。
\end{explain}

\begin{thm}
\ollabel{thm:sol-no-ls} 勒文海姆–斯科伦定理对二阶逻辑不成立：存在具有无限模型但无可数模型的语句。
\end{thm}

\begin{proof}
回想一下，
\[
\fn{Count} \ident \lexists[z][\lexists[u][\lforall[X][((X(z) \land
      \lforall[x][(X(x) \lif X(u(x)))]) \lif \lforall[x][X(x)])]]]
\]
在一个结构~$\Struct{M}$ 中为真当且仅当 $\Domain{M}$ 是可数的，因此 $\lnot\fn{Count}$
在~$\Struct{M}$ 中为真当且仅当 $\Domain{M}$ 是不可数的。存在这样的结构——取任意不可数集合作为论域，例如 $\Pow{\Nat}$ 或~$\Real$。所以 $\lnot \fn{Count}$
有无限模型但没有可数模型。
\end{proof}

\begin{thm}
存在具有可数无穷模型但无不可数模型的语句。
\end{thm}

\begin{proof}
$\fn{Count} \land \fn{Inf}$ 在 $\Nat$ 中为真，但在任何论域 $\Domain{M}$ 不可数的结构~$\Struct{M}$ 中不为真。
\end{proof}

\end{document}